% =========================================================================
% CAPÍTULO 5: ASPECTOS ADMINISTRATIVOS
% =========================================================================
\newpage

\chapter{Aspectos administrativos}

\section{Cronograma de actividades}

\noindent La Tabla~\ref{tab:cronograma_actividades} detalla la programación temporal de las actividades de la investigación.

\begin{table}[H]
\centering
\small
\begin{threeparttable}
\caption{Cronograma de actividades de la investigación}
\label{tab:cronograma_actividades}
\begin{tabularx}{\textwidth}{X c c c c c c}
\toprule
\textbf{Actividades / Fases del Proyecto} & \textbf{M1} & \textbf{M2} & \textbf{M3} & \textbf{M4} & \textbf{M5} & \textbf{M6} \\
\midrule
\textbf{Fase 1: Preparación y Marco Ético} & & & & & & \\
Revisión de literatura y estado del arte & X & & & & & \\
Levantamiento de requerimientos & X & X & & & & \\
Permisos institucionales y validación ética & X & & & & & \\
\midrule
\textbf{Fase 2: Diseño de la Solución} & & & & & & \\
Diseño de la arquitectura del sistema & & X & & & & \\
Selección y configuración del modelo base & & X & & & & \\
\midrule
\textbf{Fase 3: Recursos y Corpus Bilingüe} & & & & & & \\
Curaduría y normalización del corpus & & X & X & & & \\
Construcción del banco de consultas patrón & & & X & & & \\
\midrule
\textbf{Fase 4: Implementación del Prototipo} & & & & & & \\
Cuantización e integración RAG local & & & X & X & & \\
Adaptación del modelo (condicionada) & & & & X & & \\
Desarrollo de la interfaz de interacción & & & & X & & \\
\midrule
\textbf{Fase 5: Evaluación Experimental} & & & & & & \\
Ejecución de experimentos y captura de métricas & & & & X & X & \\
Análisis estadístico de resultados & & & & & X & \\
Revisión lingüística especializada, si se aprueba el protocolo & & & & & X & \\
\midrule
\textbf{Fase 6: Documentación y Tesis} & & & & & & \\
Redacción del informe final de tesis & & & & & X & X \\
Revisión con asesor y preparación de sustentación & & & & & & X \\
\bottomrule
\end{tabularx}
\begin{tablenotes}
\scriptsize
\item \textit{Nota.} Planificación relativa a seis meses desde la aprobación del plan; las fechas calendario se consignarán cuando se establezca el inicio del proyecto.
\end{tablenotes}
\end{threeparttable}
\end{table}

\section{Presupuesto y financiamiento}

\noindent La Tabla~\ref{tab:presupuesto_investigacion} resume el presupuesto analítico de la investigación.

\begin{table}[H]
\centering
\small
\begin{threeparttable}
\caption{Presupuesto analítico de la investigación}
\label{tab:presupuesto_investigacion}
\begin{tabularx}{\textwidth}{p{3.6cm} X c p{2.2cm}}
\toprule
\textbf{Rubro} & \textbf{Descripción / Concepto} & \textbf{Cantidad} & \textbf{Monto (S/.)} \\
\midrule
\textbf{1. Bienes y Equipos} & & & \\
Equipo de cómputo / prueba & Equipo disponible para las pruebas locales & 1 equipo & Aporte en especie \\
Almacenamiento y respaldo & Espacio de almacenamiento y respaldo de datos & Por definir & Por cotizar \\
Material bibliográfico & Fuentes de acceso abierto y materiales que requieran adquisición & Por definir & Por cotizar \\
\midrule
\textbf{2. Servicios} & & & \\
Servicios de cómputo / servidor & Procesamiento adicional, si las pruebas QLoRA lo requieren & Por definir & Por cotizar \\
Servicio de conectividad & Conectividad eventual para obtener fuentes y dependencias & Por definir & Por cotizar \\
Materiales e impresiones & Impresión y encuadernación requeridas por la escuela & Por definir & Por cotizar \\
\midrule
\textbf{3. Recursos Humanos} & & & \\
Validación lingüística & Revisión especializada, si se incorpora al protocolo & Por definir & Por cotizar \\
\midrule
\textbf{Total General} & & & \textbf{Por estimar} \\
\bottomrule
\end{tabularx}
\begin{tablenotes}
\scriptsize
\item \textit{Nota.} La fuente y distribución del financiamiento están pendientes de confirmación; el equipo disponible se registra como aporte en especie.
\end{tablenotes}
\end{threeparttable}
\end{table}

\section{Equipo de investigación}

\noindent La Tabla~\ref{tab:equipo_investigacion} presenta los integrantes del equipo de investigación y sus funciones.

\begin{table}[H]
\centering
\small
\begin{threeparttable}
\caption{Equipo de investigación}
\label{tab:equipo_investigacion}
\begin{tabularx}{\textwidth}{p{3.5cm} p{3.5cm} X}
\toprule
\textbf{Integrante} & \textbf{Rol} & \textbf{Funciones Principales} \\
\midrule
Edmil Jampier Saire Bustamante & Tesista & Diseño, implementación, registro de pruebas y redacción del plan. \\
Por designar & Asesor/a & Orientación y revisión del proyecto, una vez formalizada la designación. \\
\bottomrule
\end{tabularx}
\begin{tablenotes}
\scriptsize
\item \textit{Nota.} La asesoría y las funciones se actualizarán cuando se formalice la designación correspondiente.
\end{tablenotes}
\end{threeparttable}
\end{table}
